\subsection{P009 --- A systematic comparison and evaluation of building ontologies for deploying data-driven analytics in smart buildings}
\label{paper:P009}
\textbf{Authors:} Zhangcheng Qiang, Stuart Hands, Kerry Taylor, Subbu Sethuvenkatraman, Daniel Hugo, Pouya Ghiasnezhad Omran, Madhawa Perera, Armin Haller\par
\textbf{Major Category:} BIM and Semantic Information\par
\textbf{Secondary Category:} Reviews, Frameworks and Standards, Sensors, IoT and Data Integration\par
\textbf{Keywords:} Building ontology, Data interoperability, Brick Schema, RealEstateCore, Project Haystack, Digital Buildings, Ontology evaluation\par
\paragraph{主な主張}
Brick Schema, RealEstateCore, Project Haystack, Digital Buildingsの4 ontologyをTBoxとABoxの両面で比較した. OQuaREによるTBox評価と実建物dataによるABox評価から, 単一のuniversal ontologyは存在せず, ontology matching, alignment, harmonisationが必要と結論付けた. \cite{Qiang:2023BuildingOntologies}
\paragraph{新規性}
ontology axiomatic designの品質評価と, 実建物のHVAC, space, equipment, point dataを用いたempirical completeness/expressiveness評価を同一frameworkで実施した点に特徴がある.
\paragraph{位置付け}
smart building data interoperabilityを支えるbuilding ontology自体の比較評価が中心であるためBIM and Semantic Informationを主分類とする.
